\documentclass[11pt,a4paper]{article}

% Minimal, editable LaTeX setup
\usepackage[T1]{fontenc}
\usepackage{lmodern}
\usepackage[margin=28mm]{geometry}
\usepackage{amsmath,amssymb}
\usepackage{microtype}

% Short math shortcut (optional)
\newcommand{\RR}{\mathbb{R}}

\title{Multi-model ASPD: Formulation}
\author{}
\date{}

\begin{document}
\maketitle

\section{Setup}

\paragraph{Models and data.}
Consider $N$ pretrained language models $\mathcal M^{(1)},\ldots,
\mathcal M^{(N)}$, possibly with different architectures.
Let $\mathbf s\sim\mathcal D$ be an input text sequence drawn from a
shared distribution. Model $n$ tokenizes it as
$$
\tau^{(n)}(\mathbf s)
=(z_1^{(n)},\ldots,z_{T_n}^{(n)}),
\qquad T_n=T_n(\mathbf s).
$$
The models share the input text, but need not share tokenization or
aligned token positions.

\paragraph{Layers and indices.}
For each model $n$, select a set of layers $\mathcal L^{(n)}$ and define
$$
\sigma=(\ell,p)\in\Sigma^{(n)},
\qquad
j=(\ell,m)\in J^{(n)},
\qquad \ell\in\mathcal L^{(n)}.
$$
Here $\Sigma^{(n)}$ indexes selected activation sites (site type $p$),
while $J^{(n)}$ indexes selected weight matrices (matrix type $m$).
In the current setup, $p\in\{\mathrm{PRE},\mathrm{MID}\}$ and
$m\in\{Q,K,V,O,\mathrm{gate},\mathrm{up},\mathrm{down}\}$.

\paragraph{Activations.}
For each $\sigma\in\Sigma^{(n)}$, let
$$
H_\sigma^{(n)}(\mathbf s)
=(h_{\sigma,1}^{(n)},\ldots,h_{\sigma,T_n}^{(n)}),
\qquad h_{\sigma,t}^{(n)}\in\RR^{d_\sigma^{(n)}}.
$$
This represents the full activation sequence at site $\sigma$;
$t$ refers to a token position within model $n$.

\paragraph{Weight matrices.}
For each $j\in J^{(n)}$, define
$$
W_j^{(n)}\in
\RR^{d_{\mathrm{out},j}^{(n)}\times d_{\mathrm{in},j}^{(n)}},
\qquad
X_j^{(n)}(\mathbf s)=(x_{j,1}^{(n)},\ldots,x_{j,T_n}^{(n)}).
$$
Its linear projection at token position $t$ is
$$
x_{j,t}^{(n)}\in\RR^{d_{\mathrm{in},j}^{(n)}},
\qquad
y_{j,t}^{(n)}=W_j^{(n)}x_{j,t}^{(n)}
\in\RR^{d_{\mathrm{out},j}^{(n)}}.
$$
We omit optional biases. The matrix input $x_{j,t}^{(n)}$ need not
coincide with a selected activation $h_{\sigma,t}^{(n)}$.

% --- Continue your formulation here ---
% \section{Mechanism Definition}
\section{Architecture (current)}
This section serves as notation setup for different components and formula will be used later in formulation of post-hoc analysis. Therefore I won't mention again the overall architecture of multimodel ASPD.

Note: Current architecture has been changed, mainly with the encoder. This document also need update if future architecture changes.
\paragraph{Feature Gate}
Remind: After model processes input sequence s, we have
$$
H_{\sigma}^{(n)}(s) = (h_{\sigma, 1}^{(n)}, ..., h_{\sigma, T_n}^{(n)})
$$
In the current training setup, we choose 8 layers per model. Group of chosen layer is $L^{(n)} = \{l_1, l_2, ..., l_8\}$

At an anchor position t of model n, we concatenate activation vectors
$$
r_{p,t}^{(n)} = [h_{l_1, p, t}^{(n)}, h_{l_2, p, t}^{(n)}, ..., h_{l_8, p, t}^{(n)}] \in \RR^{8d_n}
$$
Note: index $t^{n}$ can be different between each model, depend on the way the sequence s is tokenized. This has been mentioned in setup section.

Learned projection: This step is preprocessed step before feeding to the first transformer. There are three main purpose of this:
\begin{itemize}
\item Decrease dimensionality. If model n has 8 layers, each activation vector has $d_n = 2560$ dimensions. After concatenate, $r_{p,t}^{(n)} \in \RR^{8d_n} = \RR^{20480}$ $\Rightarrow$
Large cost if we put it directly in transformer
\item Projection also allow combine information from different layers. P and v learn from input from multiple layers at once.
\item To create same input dimension between model. After projection, they will all have united, reduced dimension.
\item Limitation: Projection might lose information when reduce dimension.
\end{itemize}
$$
v_{p, t}^{(n)} = P^{(n)}_p*r^{(n)}_{p,t} + b^{(n)}_p
$$
After reshape, we got $E^{(n)}_{p,t}  =Reshape_{8512}(v_{p,t}^{(n)})$.

This is 8 learned latent tokens, not 8 token positions of input text.

First step: After go through model-specific Transformer:
$$
Z_{p,t}^{(n)} = F_p^{(n)}(E^{(n)}_{p,t})
$$
$F_p^{(n)}$ is Transformer with 2 blocks, output dimension is 8 x 512

Second step: After go through shared encoder.

Note: From this step I will write $\alpha$ as shared anchor. This is a transition from t. After projection they already have same dimension.
$$
Z_{p,\alpha}^{shared} = \operatorname{Transformer}(\operatorname{Concat}(Z_{p,\alpha}^1, SEP, Z_{p,\alpha}^2, SEP, ..., SEP, Z_{p,\alpha}^n))
$$ where SEP is separator embedding.

% Shared representation -> pooled representation
$$
\bar z_{p,\alpha}
=
\frac{1}{L}
\sum_{k=1}^{L}
Z_{p,\alpha,k}^{\mathrm{shared}}
$$

Why mean pooling?:
\begin{itemize}
\item After pooling we have $Z^{shared} \in \RR^{26\times512}$. We want to make a gate vector $\RR^C$ with C = 8192. Natural thought is get average over 26 latent tokens to form a 512 dimension vector.
\item Why not simply flatten the whole shared representation? There will be more weight for each projection head $\Rightarrow$ More expensive
\end{itemize}

% Feature logits
$$
a_{p,\alpha}
=
W_p^{\mathrm{gate}}\bar z_{p,\alpha}
+b_p^{\mathrm{gate}}
$$
This step simply upscale the vector to the dimension of C - number of components.
% Normalize and apply ReLU
$$
\tilde a_{p,\alpha,i}
=
\operatorname{ReLU}\left(
\frac{a_{p,\alpha,i}-\mu_{p,\alpha}}
{\sqrt{\sigma_{p,\alpha}^2+\epsilon}}
\right)
$$
(Refer to my latest meeting report on the solution to resolve feature dictionary collapse problems.)

% Continuous sparse gate over a batch of anchors
$$
g_p^s
=
\operatorname{BatchTopK}_K
\left(
\{\tilde a_{p,\alpha}\}_{\alpha\in\mathcal B}
\right)
$$

% Binary gate
$$
g_{p,\alpha,i}
=
1[g_{p,\alpha,i}^s>0]
$$

\paragraph{Definition of a component:} With model n, weight matrix $j \in J^{(n)}$, and component i, we define
$$
P_{j,i}^{(n)} = u^{(n)}_{j,i} (v_{j,i}^{(n)})^T
$$ where u is write vector and have dimension = layer output dimension and v is read vector and have dimension = layer input dimension.

Coefficient of contribution of write vector u:
$$
e_{j,\alpha,i}^{(n)} = g_{p(j), \alpha, i}((v_{j,i}^{(n)})^T * x_{j,\alpha}^{(n)})
$$
$$
\Delta y_{j,\alpha,i}^{(n)}
=
e_{j,\alpha,i}^{(n)}u_{j,i}^{(n)}
$$ Delta là corresponding contribution output.

This two metrics will have different use later in post-hoc analysis section, so keep in mind their difference: e is coefficient of contribution of write vector u, show how component react to input, and delta is real contribution output.

\section{Post-hoc analysis}

\paragraph {Edge score}
To begin, we care about the change of coefficient in one component downstream when we get rid of a component upstream.
Consider two components in model n, at anchor $\alpha$:
\begin{itemize}
\item $b = (j_2, i_2)$: Upstream component.
\item $a = (j_1, i_1)$: Downstream component.
\end{itemize}

For neat writing, let define:
$$
e_a = e^{(n)}_{j_1, \alpha, i_1}
\qquad
e_b = e^{(n)}_{j_2, \alpha, i_2}
$$

Consider $e_a$ as a function of $e_b$ (since $e_b$ is upstream component):
$$
e_a = f(e_b)
$$
Taylor approximation when $e_b$  change small $\Delta e_b$
$$
f(e_b + \Delta e_b) \approx f(e_b) + f'(e_b)\Delta e_b
$$

We want to ablate component b, which is bring it coefficient e to 0. We set $\Delta e_b = -e_b$
$$
\Rightarrow f(0) \approx f(e_b) - f'(e_b)e_b
$$
$$
\Leftrightarrow f'(e_b)e_b \approx f(e_b) - f(0)
$$
$$
\Leftrightarrow E_{b \rightarrow a} = \frac{\partial e_a} {\partial e_b} e_b \approx e_a(s) - e_a(s | e_b = 0)
$$

\begin{itemize}
\item If $E_{b \rightarrow a} > 0$: Ablate b will decrease $e_a$
\item  If $E_{b \rightarrow a} < 0$: Ablate b will increase $e_a$
\item  $E_{b \rightarrow a} \approx 0$: The change is insignificant.
\end{itemize}

\paragraph{Component Importance Score}
From the previous definition of component contribution defined in definition of a component section, I define normalized component contribution score

For each shared gate $(p,i)$ in model $n$,
let $J_p^{(n)}$ denote the set of weight matrices
controlled by gate stage $p$.

We define the normalized contribution score of component $i$
at weight matrix $j$ as

$$
\beta_{j,i}^{(n)}
=
\frac{
\sqrt{
\mathbb{E}_{\alpha}
\left[
\|\Delta y_{j,\alpha,i}^{(n)}\|_2^2
\mid g_{p,\alpha,i}=1
\right]
}
}{
\sigma_j^{(n)}
}
$$

where

$$
\sigma_j^{(n)}
=
\sqrt{
\mathbb{E}_{\alpha}
\left[
\left\|
y_{j,\alpha}^{(n)}
-
\mathbb{E}_{\alpha}[y_{j,\alpha}^{(n)}]
\right\|_2^2
\right]
}
$$

As an auxiliary metric, we define the within-matrix
relative contribution score:

$$
q_{j,i}^{(n)}
=
\sqrt{
\mathbb{E}_{\alpha}
\left[
\|\Delta y_{j,\alpha,i}^{(n)}\|_2^2
\right]
}
$$

$$
\rho_{j,i}^{(n)}
=
\frac{q_{j,i}^{(n)}}
{\sum_{k=1}^{C}q_{j,k}^{(n)}}
$$

Here, $\beta_{j,i}^{(n)}$ measures contribution relative
to the matrix output scale, while $\rho_{j,i}^{(n)}$
measures the component's relative dominance within
the same weight matrix. Both ratios require nonzero
denominators.

Note: Here currently there is an ambiguity on what is used on the screening step. Need decision on this in the future.

\paragraph{Parent set definition}

Consider two models $n$ and $m$.
Components $P_{j^{(n)},i}^{(n)}$ and
$P_{j^{(m)},i}^{(m)}$ are considered corresponding
if they are controlled by the same shared gate
$g_{p,\alpha,i}$, regardless of their layer
or weight matrix locations.

\textbf{Representative Component Selection.}

For each shared gate $(p,i)$ in model $n$,
we select representative components based on
their normalized contribution scores
$\beta_{j,i}^{(n)}$.

$$
\mathcal{R}_{p,i}^{(n)}
=
\left\{
(j,i):
j\in\operatorname{TopK}_{K_{\mathrm{rep}}}
\left(\beta_{\cdot,i}^{(n)}\right),
\quad
\beta_{j,i}^{(n)}>\tau_{\mathrm{rep}}
\right\}
$$

where TopK selects weight matrix indices
$j\in J_p^{(n)}$ with the largest scores.
$K_{\mathrm{rep}}$ is the maximum number of
representative components, and
$\tau_{\mathrm{rep}}\geq0$ is an optional threshold.
The within-matrix score $\rho_{j,i}^{(n)}$
can be used as an auxiliary screening metric.

\textbf{Parent Set Construction.}

Let

$$
a^{(n)}=(j_a^{(n)},i)\in\mathcal{R}_{p,i}^{(n)},
\qquad
a^{(m)}=(j_a^{(m)},i)\in\mathcal{R}_{p,i}^{(m)}
$$

be two selected corresponding downstream
components.

For each model $n$, let $\mathcal{U}^{(n)}(a)$
denote the set of candidate upstream components of $a$.

We define the parent set of component $a$ as:

$$
\operatorname{Pa}_{K_{\mathrm{par}},\tau}^{(n)}(a;\alpha)
=
\left\{
b\in\operatorname{TopK}_{K_{\mathrm{par}}}
\left(|E_{\cdot\rightarrow a}^{(n)}(\alpha)|\right)
:
|E_{b\rightarrow a}^{(n)}(\alpha)|>\tau
\right\}
$$

where TopK selects the $K_{\mathrm{par}}$
candidate upstream components in
$\mathcal{U}^{(n)}(a)$ with the largest
absolute edge scores, and $\tau\geq0$
is an optional threshold.

For two corresponding downstream components,
we obtain:

$$
\begin{aligned}
\mathcal{P}_{a,\alpha}^{(n)}
&=
\operatorname{Pa}_{K_{\mathrm{par}},\tau}^{(n)}
(a^{(n)};\alpha)
\\[8pt]
\mathcal{P}_{a,\alpha}^{(m)}
&=
\operatorname{Pa}_{K_{\mathrm{par}},\tau}^{(m)}
(a^{(m)};\alpha)
\end{aligned}
$$

\paragraph{Parent Set Comparison}

Goal: We want to compare distributions of edge scores of parent-gate sets between two models n and m.

Note: For a model, a shared gate might have multiple representative components. Therefore, we need a method to aggregate there distribution/effect score.

Since corresponding components may belong to different
weight matrices, we compare their parents using
shared-gate identities rather than matrix indices.

For an upstream component $b=(j_b,i_b)$, define

$$
\gamma(b)=(p(j_b),i_b)
$$
$p(j_b)$ because gate have type PRE/MID. Hàm gamma above là mapping from component identity to shared-gate identity.

For each downstream gate $(p,i)$, aggregate the parent
sets of its representative components:

$$
\mathcal{G}_{p,i}^{(n)}(\alpha)
=
\bigcup_{a\in\mathcal{R}_{p,i}^{(n)}}
\left\{
\gamma(b):b\in\mathcal{P}_{a,\alpha}^{(n)}
\right\}
$$

\textbf{Set Overlap.}

We measure the overlap between parent-gate sets using
the Jaccard index:

$$
J_{p,i}^{(n,m)}(\alpha)
=
\frac{
|\mathcal{G}_{p,i}^{(n)}(\alpha)
\cap\mathcal{G}_{p,i}^{(m)}(\alpha)|
}{
|\mathcal{G}_{p,i}^{(n)}(\alpha)
\cup\mathcal{G}_{p,i}^{(m)}(\alpha)|
}
$$

\textbf{Weighted Parent Distribution.}

To account for edge strength, we normalize each
edge score by the output scale of its downstream matrix:

$$
w_{b\rightarrow a}^{(n)}(\alpha)
=
\frac{
|E_{b\rightarrow a}^{(n)}(\alpha)|
\|u_{j_a,i_a}^{(n)}\|_2
}{
\sigma_{j_a}^{(n)}
}
$$

? Why not just use raw E edge score:
We are trying to measure relative impact of an upstream component on output of downstream component, instead of only the change of coefficient $e_a$

For each upstream gate identity $\xi$, define its
aggregated parent score:

$$
S_{p,i}^{(n)}(\xi;\alpha)
=
\frac{1}{|\mathcal{R}_{p,i}^{(n)}|}
\sum_{a\in\mathcal{R}_{p,i}^{(n)}}
\sum_{\substack{
b\in\mathcal{P}_{a,\alpha}^{(n)}\\
\gamma(b)=\xi
}}
w_{b\rightarrow a}^{(n)}(\alpha)
$$

The normalized parent distribution is:

$$
\pi_{p,i}^{(n)}(\xi;\alpha)
=
\frac{
S_{p,i}^{(n)}(\xi;\alpha)
}{
\sum_{\xi'}S_{p,i}^{(n)}(\xi';\alpha)
}
$$

We compare the parent distributions between models
$n$ and $m$ using Total Variation distance:

$$
D_{\mathrm{TV},p,i}^{(n,m)}(\alpha)
=
\frac12\sum_{\xi}
\left|
\pi_{p,i}^{(n)}(\xi;\alpha)
-
\pi_{p,i}^{(m)}(\xi;\alpha)
\right|
$$

The total normalized edge strength can also be retained:

$$
M_{p,i}^{(n)}(\alpha)
=
\sum_{\xi}S_{p,i}^{(n)}(\xi;\alpha)
$$

All distribution-based comparisons are defined for
anchors with nonempty representative sets and positive
total parent scores in both models.

\end{document}
